\section{Limitações e Próximos Passos}
\label{sec:limitacoes}

A revisão do modelo conceitual mostrou pontos em que o banco atual guarda
mais do que o modelo pede. Eles estão listados abaixo, com a mudança
planejada.

\begin{table}[htbp]
  \centering
  \caption{Pontos em aberto no esquema}
  \label{tab:pendencias}
  \small
  \begin{tabularx}{\textwidth}{@{}L{0.8}L{1.15}L{1.05}@{}}
    \toprule
    \poliCabTabela{Ponto} & \poliCabTabela{Problema} & \poliCabTabela{Mudança planejada} \\
    \midrule
    \campo{musics.rating} e \campo{musics.streams} & Dados derivados: repetem a média de \campo{rating} e a contagem de \campo{user\_event}, e podem divergir delas & Calcular na consulta e apagar as colunas \\
    \campo{rating} (nota de 1 a 5) & A HU-20 trocou a nota por like e dislike & Tabela de votos no lugar de \campo{rating} \\
    \campo{system\_data} & Contadores que um \campo{COUNT} já calcula; hoje só o seed escreve nela & Apagar a tabela ou trocá-la por uma \emph{view} \\
    \campo{music\_user} & As curtidas passaram para a playlist ``Músicas Curtidas'', e a API não usa mais a tabela & Apagar a tabela \\
    \campo{musics.lyrics\_id} & Coluna sem chave estrangeira: a tabela de letras ainda não existe & Criar a tabela de letras, se a funcionalidade entrar no escopo \\
    \bottomrule
  \end{tabularx}
\end{table}

Cada mudança será uma migration nova, seguindo as regras da
\cref{sec:evolucao}. Quando uma delas for aplicada, o modelo relacional da
\cref{fig:relacional} se atualiza sozinho no próximo push, e o modelo
conceitual da \cref{fig:der} já está na forma final.
